%leave all of this alone
\documentclass[10pt, fleqn]{article}
\setlength{\parindent}{.5in}
\usepackage{fontspec}
\setmainfont{Times New Roman}
%\usepackage[top=1.25in, left=0.75in, right=.75in, bottom=.75in, includefoot]{geometry}
\usepackage[top=1.25in, left=0.75in, right=.75in, bottom=.75in]{geometry}
%\usepackage{fancyhdr}
\usepackage{authblk}
\usepackage{titlesec}
\usepackage{indentfirst}
\usepackage[fleqn]{amsmath}
\usepackage{float}
\usepackage{caption} 
\usepackage{graphicx}
\usepackage{apacite}
\pagenumbering{gobble}
\captionsetup[table]{skip=8pt}
%\fancyhf{}
%\lhead[Header]{Proceedings of the Annual General Donald R. Keith Memorial Conference
%\linebreak West Point, New York, USA
%\linebreak May 04, 2023
%\linebreak \fontsize{8pt}{0}\textit{A Regional Conference of the Society for Industrial and Systems Engineering}}
%\renewcommand{\headrulewidth}{0pt}
\setlength{\headheight}{0pt}
\renewcommand\Authands{, and }
\renewcommand{\Authfont}{\bfseries\fontsize{12pt}{0}}
\renewcommand{\Affilfont}{\normalfont\vspace{0.5em}\normalsize}
\date{}
\titleformat{\section}
  {\centering\fontsize{11pt}{0}\bfseries}
  {\thesection.}
  {.5em}
  {}
\titleformat{\subsection}
  {\normalfont\fontsize{11pt}{0}\bfseries}
  {\thesubsection.}
  {.5em}
  {}
\titleformat{\subsubsection}
  {\normalfont\fontsize{11pt}{0}\bfseries}
  {\thesubsubsection.}
  {.5em}
  {}
\titlespacing*{\subsubsection}{\parindent}{1em}{0ex}


%put your title here
% \title{\vspace{-1cm}\fontsize{14pt}{0}\bfseries\textbf{
%     Paper Title
% }}
\title{\vspace{-14pt}\fontsize{14pt}{0}\bfseries\textbf{
    Paper Title 
}\vspace{-0.5em}}

%put authors and affiliations in like this
%give the corresponding author a * and list their email below in the affiliations
\author[1]{(First Name Last Name) Author}
\author[2]{Author}
\author[3]{Author}
\affil[1]{Department of Systems Engineering, United States Military Academy, West Point, New York  10996}
\affil[2]{Author2 Affiliation (if different from Author$^1$}
\affil[3]{Author3 Affiliation (if different from Authors$^{1\&2}$)}
\affil[ ]{Corresponding author's Email: xyz@abc.com (use an email that will not expire after graduation)}


%leave all of this alone
\begin{document}
%\pagestyle{fancy}
%\fancyfoot{}
%\renewcommand{\footrulewidth}{0pt}
%\fancyfoot[C]{\thepage}
%\fancyfoot[L]{ISBN:}
\maketitle
%\thispagestyle{fancy}

%author note - leave the next line alone
\vspace{-1cm}\noindent\textbf{Author Note: }
%add your author note text here...
A brief biographical sketch of the contributing authors goes here. It may include acknowledgment of funding sources, expressions of gratitude to research assistants, and contact information for the author who will handle requests. Projects written by cadets and ROTC students should include a disclosure statement similar to, “The views expressed herein are those of the authors and do not reflect the position of the United States Military Academy, the Department of the Army, or the Department of Defense.”


%abstract - leave the next line alone
\vspace{1em}\noindent\textbf{Abstract: }
%add your abstract text here...
This article illustrates preparation of your conference paper using Overleaf. All papers should be submitted both in PDF format and as a zip file with all of your project files.  You can download a PDF by clicking "compile" in the Overleaf editor and then "download" to export the result.  You can download the project as a zip file by navigating to your list of Overleaf projects and clicking "download". \textbf{Papers are limited to 6 pages and abstracts are limited to 150 words.} Please modify this template \textbf{without modifying the formatting options}. 

%Add your keywords here...
\vspace{1em}\noindent\textit{Keywords}: Keyword1, Keyword2, Keyword3

%Add your paper content here - pay attention to the sections, subsections, and subsubsections.
%Otherwise, it should be easy to add your content.
\section{Heading Level 1}

This is the beginning of the body of your paper. The rest of the text in this template provides hints about properly generating all the parts of your APA-formatted paper.  Note that all of the formatting is done for you so you should be able to paste your text into this document without worrying about things like spacing and indentation.

Sections can be further divided into subsections with headings.  


\subsection{Heading Level 2}

The first  paragraph after any section title will be indented as here with one line space before the paragraph.

The second and all subsequent paragraph will also be indented (0.5 inches) as here with no line spacing before the lines


\subsubsection{Heading Level 3}
The first paragraph after any section title will be indented as here with no line space before the paragraph.

The second and all subsequent paragraphs will also be indented as here with no line spacing before the line. Avoid using more than three levels of headings.

\section{Equations}

Enter equations as shown below.  This will automatically generate a number for the equation and format it correctly.  Use a "LaTeX Formula Editor" to help build these if you are unfamiliar with LaTeX formula notation.

%note how we add equations
\begin{equation}
f(t) = \int_{0_+}^{t}F(t)dt+\frac{dg(t)}{dt}
\end{equation}


\section{Figures and Tables}
Below is an example of a LaTex table.  Feel free to use other styling in your tables not shown here, but please ensure that the font is 9pt or larger.  If you want the table to render exactly where you put it in your .tex file, you will need to add [H] as shown below.  As with formulas, you may find online "LaTeX Table Editors" to be helpful.

%note how we add tables
\begin{table}[H]
\centering
\caption{Correlations Among and Descriptive Statistics For Key Study Variables}
\begin{tabular}{llllllll} 
\hline\hline
             & M (SD)        & Sex & Age & Income & Educ.                  & Relig.                  & Dist. Intol.             \\ 
\hline
Sex          & 1.53 (.50)    &     & .07 & -.09   & .02                    & .14                     & .06                      \\
Age          & 31.88 (10.29) &     &     & .08    & .19\textsuperscript{*} & .20\textsuperscript{*}  & .01                      \\
Income       & 2.60 (1.57)   &     &     &        & .04                    & -.14                    & -.09                     \\
Education    & 3.44 (1.57)   &     &     &        &                        & -.29\textsuperscript{*} & -.06                     \\
Relig.       & 1.21 (.30)    &     &     &        &                        &                         & -.19\textsuperscript{*}  \\
Dist. Intol. & 3.75 (1.19)   &     &     &        &                        &                         &                          \\
\hline
\end{tabular}
\end{table}

All figures should be positioned at the top of the page when possible. All figures will be numbered consecutively and captioned; the caption should be centered under the figure as shown in Figure 1. All text within the figure should be no smaller than 9pt.  As with tables, you will need to specify [H] as shown below to ensure the figure appears in the expected location relative to your text.
c%note how we add figures
\begin{figure}[H]
    \begin{center}
        \includegraphics{figure_1}
        \caption{Title of Figure}
    \end{center}
\end{figure}


\section{Citations and References}

Citations and references should be generated using Bibtex.  Add your references to the refs.bib file provided along with this template.  Bibtex citations are typically provided by academic journals and can also be pulled from Google Scholar.  

Within the text, references should be cited giving the last name of the author(s) and the year of publication of the reference. You can generate these citations using Bibtex as well - for example: \citeA{allen1974vanishing} or \cite{allen1974vanishing}. 

Your references section will be automatically generated at the end of your paper in the correct format.

%leave this alone
\bibliographystyle{apacite}
%the next line adds all of your references even if you didn't cite them in the paper.
%to override that, you can just delete the next line
\nocite{*}
\bibliography{refs}

\end{document}


